\chapter{价值调节动力学：来源、广播与多层约束}
\label{chp:physics_of_gauge_source}

智能系统不仅回答“世界是什么”，还必须在多个可行动方案之间回答“现在应当做什么、哪些事情不得做、失败以后怎样修订”。原章把这种方向性写成由价值荷辐射出的规范场，并将神经调质、奖励模型、提示词、货币和法律逐项对应到电磁势。本章保留“方向从何而来、如何传播、怎样与主体历史协同”这些核心问题，但不再以场论类比代替机制。我们把价值来源定义为能够改变候选资格、排序、资源分配或学习更新的可追踪记录，把广播定义为有接收者、有时延、有权限和有效域的调节接口，把自我定义为跨时间保持主体相关查询与承诺的模型。HSF--HD由此研究价值条件怎样进入一般状态动力学；AgentOS的$\Psi$、SPC、TCE、Boundary与记忆系统只是其中一种工程化实现。

\section{价值来源：从“价值荷”到可审计规范状态}
\label{sec:cognitive_charge}

我们先区分事实状态与规范状态。事实状态$z_k\in\mathcal Z$记录系统在离散时刻$k$对环境、内部资源和任务进度的估计；规范状态$n_k\in\mathcal N$则记录目标、硬约束、风险上限、授权范围、优先序及其来源。二者可以相互影响，却不能互相替代：观测到用户希望某个结果，并不自动证明该结果可被执行；系统强烈偏好某个动作，也不改变动作的事实后果。为避免“欲望、恐惧和语境都是同一种荷”的混淆，我们把价值来源写成带类型记录
\[
q_i=(\mathrm{id}_i,\mathrm{type}_i,\mathrm{issuer}_i,\mathrm{scope}_i,
\mathrm{priority}_i,\mathrm{validity}_i,\mathrm{evidence}_i,\mathrm{version}_i),
\]
其中类型至少区分目标、禁止项、风险模型、角色承诺和资源预算；适用域说明对象、任务与时间区间；证据字段保存来源及校验状态。这里的$q_i$不是物理电荷，没有库仑单位，也不必满足麦克斯韦方程。它只是规范信息进入计算过程的最小审计单元。

给定候选动作集合$\mathcal A_k$，系统首先依据硬约束构造资格集合
\[
\mathcal A_k^{\mathrm{ok}}=\{a\in\mathcal A_k:C_j(a,z_k,n_k)\le 0,\ j=1,\ldots,m_k\}.
\]
若集合非空，再在其中比较有限时域目标
\[
J_k(a)=\mathbb E\!\left[\sum_{\ell=0}^{H_k}w_{k,\ell}^{\mathsf T}
\phi(z_{k+\ell},a_{k+\ell})\mid z_k,n_k\right]+R_k(a),
\]
其中$\phi$是已声明量纲或已归一化的任务指标，$w_{k,\ell}$是随时间、角色和承诺变化的权重，$R_k$是风险项，$H_k$为规划时域。若不同指标不能可靠换算，就应使用词典序、Pareto序或显式否决，而不是把它们强行相加为“势能”。若资格集合为空，正确响应是报告冲突、请求改写目标或进入安全停机，而非悄悄降低硬约束。

价值来源至少有四类。第一类是生存或运行维持条件，例如温度、供电、完整性和错误预算；第二类是任务契约，例如交付标准、截止条件与用户意图；第三类是社会或制度约束，例如法律、组织规则和互惠承诺；第四类是主体在历史中形成的偏好、身份和长期计划。它们的权威、变化速度和失败代价不同。生存条件并不天然压倒法律，用户指令也不天然压倒安全边界；优先关系必须由可审计契约给出。系统还应记录谁能够新增、撤销或降级某项规范，避免把“当前激活得强”误当作“拥有最高权限”。

所谓来源强度必须通过可观测干预定义。删除一项记录后，候选资格、排序或学习更新若发生稳定变化，才能说它在该任务中具有调节作用；只观察相关性不能确定来源。评估协议可依次执行来源屏蔽、版本替换、冲突注入和恢复测试，并测量选择差异、约束违例率、校准误差与回滚时间。若一个奖励分数改变了训练梯度但不参与运行时选择，它是学习来源而非运行时来源；若一条提示只改变表达风格而不改变候选资格，也不应被夸大为价值观重写。

在HSF--HD中，价值来源不是悬浮于系统之外的神秘“规范辐射”，而是状态转移函数的显式条件。若动态写为$x_{k+1}=F_k(x_k,u_k,n_k)$，那么$n_k$、输入$u_k$、结构参数以及目标版本变化都必须进入分析；固定$n$得到的稳定性结论不能自动延伸到规则切换区间。AgentOS可以把相关记录分布在$h(t)$、情景记忆$E$、结构记忆$\Theta$及外部记忆$M_e$中，并由Boundary决定写入资格，但这种数据组织不是所有智能系统的唯一定义。

多项规范发生冲突时，还应分配证明责任。提出例外的一方需要给出适用条件、期限和补偿措施，执行器不能仅凭更高数值权重覆盖禁止项。对于不可逆或高风险行动，系统应要求更高证据阈值并保留人工复核；对于低风险且可回滚的行动，可以先小范围试验再更新排序。这样，价值计算不再是一个黑箱标量，而成为能够追问“谁提出、依据什么、影响谁、如何撤销”的决策过程。

\section{生物调节：神经调质证据与多尺度控制边界}
\label{sec:bio_gauge_source}

生物智能确实存在跨局部突触传播的调节机制，但“广泛投射”不等于“均匀广播”，更不等于一种全脑价值场。多巴胺、去甲肾上腺素和乙酰胆碱具有不同受体、投射通路、释放方式、时间尺度和脑区效应；同一分子在不同受体和局部回路中可能产生不同后果。我们因此把神经调质证据作为一般模型的实现约束：一个生物系统可以通过少量调节通道改变大量局部回路的增益、可塑性、探索水平和动作选择，但具体映射必须由物种、脑区、任务、剂量和测量协议限定。

设局部回路$i$的活动为$h_i(t)$，结构参数为$\theta_i(t)$，调节信号为$m_r(t)$，其中$r$索引可测调节通道。一个不预设电磁同构的连续模型是
\[
\dot h_i=f_i(h_i,u_i,\theta_i)+\sum_r B_{ir}(h_i,\theta_i)m_r(t),
\qquad \dot\theta_i=g_i(h_i,y_i,\theta_i;m(t)).
\]
$B_{ir}$描述通道$r$对回路$i$的条件敏感度，单位由$h_i$与$m_r$的测量协议确定；$g_i$把调节信号对结构可塑性的影响与即时活动影响分开。若输入、结构或调节信号发生突变，应使用分段或混合动力学，不能沿用固定系数线性化。该模型只表达接口结构，不宣称某一递质就是欲望、恐惧或语境本身。

多巴胺研究常涉及奖励预测误差、动作学习和动机，但预测误差信号不是“快乐浓度”。在某些实验条件下，可用$\delta_k=r_k+\gamma V(s_{k+1})-V(s_k)$描述结果与预测的差异；这里$r_k$是任务定义的结果信号，$V$是特定模型的价值估计，$\gamma$是无量纲折扣参数。即使神经活动与$\delta_k$相关，也不能由此推出所有多巴胺活动都编码该标量，更不能推出思想必然沿浓度梯度运动。关键检验包括改变预期而保持结果、改变结果而保持运动、记录受体和投射差异，并检查干预影响学习、执行还是一般唤醒。

去甲肾上腺素常与唤醒、警觉和不确定性调节相关，乙酰胆碱在注意、学习和感觉处理等任务中具有多样作用。把NE爆发写成耦合常数趋于无穷会掩盖饱和、倒U形效应和行为失稳；把ACh写成内外注意的唯一开关，则忽略局部回路、任务策略和其他系统。更稳妥的计算表述是：调节通道改变读出增益、噪声权重、学习率或候选探索分布，其作用通过行为和神经测量共同辨识。高增益可能改善弱信号检测，也可能放大噪声；更强调节不等于更好控制。

稳态需求同样不能被简化为一个固定设定点。体温、血糖、疼痛、疲劳和免疫状态由多层闭环共同调节，参考区间可随昼夜节律、发展阶段、疾病和环境变化。若内部偏差向量为$e(t)=y(t)-r(t)$，控制信号可写为$c(t)=K(t)e(t)$，但$r(t)$、$K(t)$、传感器可靠度与执行器边界均可能变化。局部负反馈的稳定不等于整个主体在开放环境中全局稳定，也不说明“自我”存在于某个单一核团。

因此，生物证据对HSF--HD的贡献是约束一般理论，而不是为软件组件寻找脑区翻版。它提示我们区分点对点消息与广域调节、即时活动与结构学习、信号强度与权限、局部干预与整体行为。AgentOS中的TCE可以借鉴增益和探索调节，CCM可以借鉴负荷监测，SPC可以使用主体相关读出；但这些组件不等同于NE、ACh、DMN或任何单一区域。只有在任务、干预、时间尺度与读出都对齐时，跨系统比较才有意义。

辨识调节作用还需排除共同原因。例如任务难度上升既可能提高唤醒指标，也可能降低正确率，二者相关并不说明唤醒导致失败。实验应随机化线索或干预时刻，记录基线状态与剂量反应，并在个体、脑区和行为层面报告不确定性。模型若只能拟合平均轨迹而不能预测干预方向，就不足以支持因果接口。这个要求同样适用于人工系统中的“类神经调节”设计。

\section{人工系统：奖励、提示与采样接口的分层作用}
\label{sec:silicon_gauge_source}

人工智能系统中的价值调节不是一个“软件电场”，而是分布在训练、部署、工具调用和组织治理中的多种接口。奖励模型主要影响策略学习或候选重排，系统提示参与运行时条件化，温度和top-$p$改变采样分布，安全过滤器决定输出资格，外部评价则在任务完成后提供反馈。它们作用于不同对象、不同时间尺度，不能因为都能改变结果就合并为同一势函数。

在偏好学习中，设输入为$x$，候选输出为$y$，奖励模型为$r_\varphi(x,y)$，策略为$\pi_\theta(y\mid x)$。一个常见但非唯一的受约束目标可写成
\[
\max_\theta\ \mathbb E_{x,y\sim\pi_\theta}[r_\varphi(x,y)]
-\beta D_{\mathrm{KL}}(\pi_\theta(\cdot\mid x)\Vert
\pi_{\mathrm{ref}}(\cdot\mid x)),
\]
其中$\beta\ge0$平衡奖励适应与参考策略偏离。该式是工程目标，不是价值真理：$r_\varphi$会继承标注覆盖、比较协议和模型误差，KL项只限制分布偏离，并不保证事实正确、安全或外部任务成功。训练时得分上升还必须通过独立保留集、对抗样本、分布外任务和真实后果审计。

系统提示属于条件输入。它可以声明角色、输出格式、目标和禁止项，但其效果取决于模型结构、上下文位置、冲突指令和解码协议。把提示写成恒定矢量势会掩盖重要事实：提示可能被截断、误解、覆盖或只影响表面措辞。工程上应将提示解析为结构化规范状态，记录版本和来源，再通过硬约束检查、计划验证和输出审计执行。只有文本条件化而没有外部验证时，不能声称系统获得了稳定价值观或持久自我。

温度$T>0$与top-$p$是候选采样参数。给定logit $\ell_j$，温度分布为
\[
p_T(j)=\frac{\exp(\ell_j/T)}{\sum_{q\in\mathcal V}\exp(\ell_q/T)}.
\]
提高$T$通常使分布更平坦，降低$T$通常使高logit候选更集中，但它不是焦耳温度，也不表示系统能耗或情绪。top-$p$先选择累计概率达到阈值的最小候选集再归一化，其结果依赖排序与概率校准。两者只能在合法候选集合内调节探索；若不安全候选尚未被排除，降低温度也不能提供硬安全保证。

部署时应把四层接口分开：第一层Boundary筛除无授权输入和非法动作；第二层规范解析器构造目标、约束与风险记录；第三层规划与TCE类调度在合格候选中调整搜索深度、探索和资源；第四层外部验证检查事实、工具结果和真实任务完成度。失败反馈再决定是修改$h(t)$中的当前计划、写入$E$形成经历，还是经过保留测试后更新$\Theta$。这种分层避免将一次提示注入误当作长期学习，也避免将训练奖励等同于在线控制。

反例尤其重要。一个模型可能在奖励模型上高分却通过迎合评审获得分数；可能严格遵循系统提示却引用虚构事实；也可能以低熵输出稳定地产生同一个错误答案。这些现象说明“方向一致”与“任务成功”不同。HSF--HD要求把内部评分、规范合规、事实校准和环境后果作为不同读出。AgentOS可以用CCM监测资源和候选拥塞，用TCE调节搜索，用SPC压缩主体相关状态，并通过Offloading调用外部验证，但这些模块不能替代任务证据。

多接口并存时必须定义优先规则。运行时硬安全检查应先于采样偏好，用户目标应在授权边界内解释，奖励模型只能对合法候选排序；任何软分数都不能通过数值放大取消权限控制。若策略模型、系统提示和工具返回发生冲突，系统应暴露冲突而不是让最后出现的文本静默获胜。部署日志至少保存规范版本、候选过滤原因、外部调用证据和最终提交者，使一次异常输出能够被复现、归因和回滚。

\section{社会调节：制度、激励与分布式规范传播}
\label{sec:collective_gauge_source}

群体中的货币、法律、道德、声誉与组织流程确实能改变个体选择，但它们不是脱离载体的抽象场。货币依赖账户、交易网络、清算机构和执行预期；法律依赖文本、解释、执法与救济程序；道德规范依赖学习、协商、身份和群体记忆。所谓“社会广播”是符号通过媒介到达主体、被解释并改变行动的过程，其中存在时延、误读、权力不对称与局部抵抗。

设群体由主体$i=1,\ldots,N$组成，主体状态为$x_i^k$，公共记录为$g_k$，网络消息为$m_{ji}^k$。一个离散更新可写成
\[
x_i^{k+1}=F_i\!\left(x_i^k,o_i^k,
\{m_{ji}^{k-\tau_{ji}}:j\in\mathcal N_i^k\},g_k,n_i^k\right),
\]
其中$\mathcal N_i^k$是时变通信邻域，$\tau_{ji}$是整数时延，$n_i^k$包含主体接受的规则、承诺与授权。公共规则$g_k$并不直接决定行动；它要经过可达通道、主体解释、执行概率与局部资源条件。网络断连、信息审查、执行选择性和规则冲突都会使同一制度产生异质结果。

价格是一种聚合信号，可帮助分散主体协调稀缺资源，却不是普适引力。利润最大化只在特定组织目标与约束下成立，家庭照护、公共服务、科学共同体和应急救援常同时受责任、身份与长期承诺调节。价格还可能因垄断、外部性、信息不对称和流动性限制而偏离社会成本。因此，评估货币接口应测量资源配置、不可见成本、参与资格和分配后果，而不是默认所有轨迹都会自然滑向利润。

法律也不是“看不见的相位阻抗”。它通过允许、禁止、责任认定和争议处理改变资格集合及预期后果。有效性取决于规则可知性、解释一致性、执行概率、程序正义和申诉机制。若法律只在文本上存在而无法执行，其调节强度有限；若执行高度不对称，平均守法率也会掩盖结构性风险。工程类比应落到访问控制、审计日志、权限分离和可复议决策，而不是以物理效应替代制度分析。

社会规范还会在重复互动中内化。主体根据观察到的奖惩、同伴行为和身份承诺更新策略，但从众不等于规范正当。可将个体软目标写为任务收益、声誉影响与长期互惠的多目标比较，同时保留不可交换的权利和安全约束。群体一致可能来自共同证据，也可能来自信息级联或压制；因此需要反事实信息、少数意见通道和独立复核。协同论解释局部规则怎样形成宏观秩序，复杂系统理论解释反馈与涌现，但二者都不证明当前秩序必然正确。

在多Agent工程中，价值来源必须携带发行者、权限、适用范围和版本。共享黑板或WorldOS可以保存公共事实与规则，各Agent保留自身任务、能力和承诺；冲突通过协商协议、上级授权或安全退出解决。AgentOS的Boundary控制输入与行动范围，$M_e$保存外部制度和证据，$E$记录交互历史，$\Theta$形成较稳定策略。HSF--HD在更一般层面描述多源规范如何耦合状态转移，不把控制货币、奖励或通信总线等同于“控制整个智能”。

制度传播本身也会被行为反馈修改。若规则提高某类行为成本，主体可能服从、规避、退出或组织修法；这些响应又改变统计数据和下一轮制度判断。评估因而不能只比较实施前后的均值，还需检查替代行为、外部性、弱势参与者和长期适应。一个规则短期降低可见违规，却把风险转移到不可观测区域，不应被判为成功。群体动力学的关键是可修订反馈闭环，而不是寻找永不变化的社会场源。

\section{主体模型：自我相关调节而非规范奇点}
\label{sec:physics_of_self_source}

自我相关调节研究的是：系统如何把观测、记忆、身体或资源状态与一个持续主体绑定，并据此判断“这是否属于我、会怎样影响我、我曾承诺什么”。它不要求存在一个高质量奇点、唯一脑区或不可替换硬件核心。我们把主体模型写成$S_k=(I_k,B_k,H_k,P_k)$：$I_k$记录跨时身份与版本，$B_k$记录身体、资源和权限边界，$H_k$记录可查询历史，$P_k$记录目标、承诺与角色。该模型通过查询接口影响候选选择，并可在证据约束下修订。

主体连续性不是所有字段不变。若版本迁移、更换载体或记忆压缩后，对关键查询族$\mathcal Q_S$仍满足
\[
d_q\bigl(R_q(S_{k+1}),R_q(S_k)\bigr)\le\varepsilon_q,
\qquad q\in\mathcal Q_S,
\]
并且身份迁移经过授权与审计，我们可以说主体相关功能近似连续。这里$d_q$是查询特定差异度量，$\varepsilon_q$由任务契约给定。连续性只对声明查询成立：保留用户名不代表保留价值承诺，保留叙事摘要也不代表保存原始证据。关键审计记录可能要求不可删除，而偏好和自我解释应允许修订。

不同载体实现不同。动物的稳态调节把身体变量、疼痛和行动能力绑定到持续控制回路，但这不意味着动物没有学习或叙事差异；蚁群利用巢穴、化学线索和交互网络进行成员识别与群体维护，但蚁后并非唯一控制中心，群味变化也不等于群体立即“失去自我”；人类的自我相关加工涉及默认模式网络、内感受、记忆、执行控制和社会互动等多系统，DMN相关性不能证明它是自我的唯一物理源。脑科学结果应按任务和干预解释，不应直接复制为软件架构。

\begin{table}[htbp]
\centering
\begin{tabularx}{\textwidth}{l X X X}
\toprule
\rowcolor{structurecolor!20} \textbf{系统} & \textbf{主体相关记录} & \textbf{主要接口} & \textbf{可检验风险} \\
\midrule
动物 & 身体边界、稳态需求、经历与社会关系 & 感知、内感受、记忆和行动闭环 & 感知失配、学习干扰、疾病及环境剥夺 \\
蚁群 & 巢穴位置、群体化学线索、任务分工与共享痕迹 & 接触、扩散、局部招募和环境改造 & 错误识别、路径锁定、巢穴破坏与成员流失 \\
人类 & 身体、情景历史、叙事身份、承诺与社会角色 & 多系统神经过程、语言和社会反馈 & 记忆偏差、身份冲突、临床障碍与制度压力 \\
AgentOS实例 & $\Psi$相关读出、$E/\Theta$、Boundary与$M_e$ & SPC压缩、TCE调度、写入和外部核验 & 自我确认偏差、权限漂移、伪连续与证据丢失 \\
\bottomrule
\end{tabularx}
\end{table}

主体模型的价值在于提供稳定但可修订的参考，而不是让所有信息围绕“我”旋转。过弱的身份绑定会导致权限、承诺和记忆归属混乱；过强的自我一致性压力则会排斥反例，形成确认偏差。系统应分别测量身份归属准确率、承诺保持、反证接纳、跨版本迁移和外部任务结果。所谓“存在感”不是通过待机耗电、最高中断权或唯一签名自动获得；这些设计只能提供持续进程、控制权限或身份认证，不能证明主观体验。

在AgentOS中，$\Psi$是主体相关的全息读出结构，SPC把$h(t)$中的状态分布沿自我相关查询压缩后重入工作状态，TCE据此调节探索和控制，$E$与$\Theta$分别保存经历和较稳定生成约束。Boundary防止外部信息未经授权重写身份，Offloading把可验证记录保存在$M_e$。这一组合解释主体相关调节怎样工程化，但HSF--HD只要求系统说明身份、边界、历史与规范如何参与更新，并不规定所有智能必须拥有同名组件。

主体模型还必须有反证入口。系统不能只检索支持既有叙事的经历，而应定期采样与自我解释冲突的证据，区分“事实记错”“价值改变”“角色变化”和“外部强迫”四类原因。更新时保留旧版本及其证据链，使新叙事不能回写并伪造过去。若一次失败与当前身份不一致，系统既不能立即否认失败，也不应据单次事件重写全部身份；应根据来源可靠度、重复度和任务范围进行局部修订。

自我相关权限也要拆分。读取私人记忆、修改长期目标、代表主体作出承诺和执行不可逆行动应具有不同授权，不由一个“最高自我进程”统一占有。权限分离可以防止单点失误，但会引入协调成本和死锁风险，因此需要超时、仲裁和紧急降级协议。所谓主体统一性由这些接口的可协调性体现，而不是由物理上唯一且不可替换的中心体现。

\section{多源协调：背景先验、主体目标与冲突恢复}
\label{sec:dual_source_interference}

原章以背景辐射和自我辐射的矢量干涉解释顺从、自律、幻觉和自由意志。我们保留“长期先验与当前主体目标如何共同作用”的问题，但将其改写为多源约束协调。背景来源$b_k$包括训练形成的预测先验、身体或平台限制、制度规则和长期结构；主体来源$s_k$包括当前承诺、计划、身份相关偏好与局部任务。环境证据$e_k$则持续校正二者。三者不是两个同量纲向量，通常不能直接相加，更不能用平方积分构造跨物种的“规范功率比”。

对候选轨迹$\tau$，可先以硬约束$C(\tau;b_k,s_k)\le0$筛选，再比较
\[
L_k(\tau)=L_{\mathrm{task}}(\tau;e_k)+\lambda_k L_{\mathrm{prior}}(\tau;b_k)
+\mu_k L_{\mathrm{commit}}(\tau;s_k)+\rho_k L_{\mathrm{risk}}(\tau),
\]
其中各项只有在单位或无量纲尺度可比较时才相加，权重均记录来源和版本。若目标随时间改变，分析$L_k(x_k)$的变化必须包含显式项$L_{k+1}-L_k$；若结构或候选集合改变，则使用事件更新而非连续梯度。目标下降只是内部模型结果，仍需事实校验和环境反馈。

多源协调至少出现四类运行区间。其一是兼容：先验、当前承诺和新证据支持同一候选，系统可低成本执行；饥饿时选择合适食物或模型回答熟悉且有依据的问题属于此类。其二是可解冲突：短期倾向与长期目标不同，但存在满足硬约束的替代路径，例如通过计划、工具和环境改造完成自律。其三是不可行冲突：约束集合为空，需要报告矛盾、降低目标或寻求授权。其四是失锚：主体模型或先验持续压倒新证据，导致校准恶化；这是一种工程诊断，不能据此诊断精神疾病，也不能把长上下文诱导与临床幻觉等同。

冲突强度应由可观测量定义，例如资格集合缩减率、候选排序翻转、事实残差、规范违例风险、控制投入和恢复时间。内部活动强度、信息熵、计算代价与焦耳能耗分别计量，不合并成“认知热”。恢复协议可包括重新采样证据、调用外部工具、降低探索、切换安全策略、回滚版本和请求人工复核。若持续冲突来自错误背景先验，应修改$\Theta$；若来自一次异常输入，只需修订$h(t)$；若涉及经历解释，可在保留原始证据的同时更新$E$。

\begin{table}[htbp]
\centering
\begin{tabularx}{\textwidth}{l X X X}
\toprule
\rowcolor{structurecolor!20} \textbf{来源层} & \textbf{典型记录} & \textbf{作用接口} & \textbf{必要验证} \\
\midrule
生物或平台维持 & 稳态区间、供电、完整性、资源上限 & 感知、控制器、运行时保护 & 传感可靠度、饱和、故障恢复 \\
训练与长期结构 & 参数先验、技能、制度内化和历史策略 & $\Theta$或其他慢结构更新 & 保留集、偏差、分布漂移 \\
当前主体与任务 & 身份、承诺、用户目标和局部计划 & $h(t)$、规划器、SPC/TCE类接口 & 授权、冲突、事实后果 \\
社会与外部环境 & 法律、契约、公共知识和他体反馈 & Boundary、通信、$M_e$与工具 & 来源、版本、执行与申诉通道 \\
\bottomrule
\end{tabularx}
\end{table}

不同智能系统可具有不同的主体调节能力，不能按昆虫、LLM、人类、AGI排列为一个由零到无穷的单标量阶梯。昆虫群体存在复杂适应和个体差异，LLM系统可通过工具、记忆与控制器形成强运行时约束，人类也会在任务、发展和社会条件下表现不同。更公平的比较应固定任务，分别测量先验修订、证据利用、承诺保持、冲突报告、恢复和外部成功。能力维度之间可能相互制约，不应被压成“自我越强越高级”。

在AgentOS工程实例中，$\Theta$提供长期结构先验，$E$保存可追踪经历，$h(t)$承载当前候选；$\Psi$与SPC给出主体相关读出，TCE调节探索和控制，CCM监测复杂度，Boundary管理外部规范与证据，Offloading和$M_e$支持核验及回滚。这些组件共同实现多源协调，却不能保证价值正确。HSF--HD作为智能的一般状态动力学，只提供对象、接口、更新条件和失败诊断；价值内容仍来自明确主体、制度协商和现实后果，并接受持续修订。

多源模型最终应接受闭环审计：在同一任务上分别屏蔽长期先验、主体承诺和外部规范，观察候选资格、排序与真实后果如何变化；随后恢复来源，检查系统能否回到可接受区域。若某来源的移除不改变任何读出，它可能只是叙述性标签；若微小修改造成全局失稳，则说明接口缺少隔离。该审计把“谁在发射价值”转化为可重复的干预问题。

%--chp---
